\documentclass[11pt]{article}

% preprint = non-anonymous, page numbers shown (suitable for course submission)
\usepackage[preprint]{acl}

\usepackage{latexsym}
\usepackage[scheme=plain,UTF8]{ctex}
\usepackage{microtype}
\usepackage{inconsolata}
\usepackage{graphicx}
\usepackage{booktabs}
\usepackage{enumitem}
\setlength\titlebox{5cm}

\title{Speeding Up Autoregressive LLM Generation via Speculative Decoding:\\[1mm]
Principle, Self-Implementation, and Evaluation}

\author{
  Zou Bowen (邹博闻) \and Cui Di (崔迪) \\
  College of Computer Science, Nankai University\\
  Team: 创造性AI \qquad \texttt{2312251@mail.nankai.edu.cn} \qquad \texttt{19837008208@163.com}
}

\date{October 12, 2026}

\begin{document}
\maketitle

\section{Background and Motivation}
Autoregressive language models generate text one token at a time: each new token
requires a full forward pass over the entire token history. For large models this
step is \emph{memory-bandwidth bound}---reading the weight matrices from memory each
step dominates compute---which limits throughput to tens of tokens per second on a
single accelerator.

\textbf{Speculative decoding} \citep{leviathan2023fast, chen2023accelerating}
reduces the number of expensive target-model forward passes without changing the
output distribution. A cheap \emph{draft model} proposes $K$ speculative tokens;
the target model verifies them in a \emph{single} parallel forward pass, and
tokens are accepted probabilistically so that the final distribution is identical
to vanilla autoregressive sampling. Because small models agree with large ones on
most frequent short tokens---a statistical property of natural-language token
distributions---most drafts are accepted, yielding often 2--3$\times$ speedups
\emph{at zero quality cost}.

Rather than merely \emph{using} an existing engine (explicitly forbidden by the
project rules), this project implements the speculative-decoding pipeline
\emph{from scratch} over an open-source LLM, rigorously verifies the ``lossless''
claim, and delivers a runnable interactive demo.

\section{Problem Definition and Research Questions}
Given a target LM $p$, a draft LM $q$, a prompt, and draft length $K$, we implement
a speculative decoding loop whose generated distribution equals vanilla
autoregressive decoding. We ask:

\begin{enumerate}[label=RQ\arabic*.]
\item Can a self-implemented engine reproduce the claimed lossless speedup (exact
      output preservation) in both greedy and sampling modes?
\item How do drafting strategies (small LM vs.~n-gram prompt lookup) and
      hyper-parameters ($K$, temperature) affect throughput, acceptance rate, and
      oracle-like quality?
\item How do gains scale with the draft/target model pair, sequence length, and
      serving batch size?
\end{enumerate}

\section{Technical Approach and System Design}
\paragraph{Core modules (all self-implemented).}
\begin{enumerate}[label=\arabic*.]
\item \textbf{Decoder engine}: KV-cache autoregressive generation, greedy and
      sampling modes (the baseline).
\item \textbf{Draft module}: $K$-step continuous sampling from a small draft LM;
      plus an n-gram prompt-lookup variant that proposes tokens copied from the
      prompt (no second model needed) \citep{nvidia2023promptlookup}.
\item \textbf{Validator}: batch forward pass of the target model over the $K$
      proposed positions; per-position acceptance with probability
      $\min\!\bigl(1,\; p(\cdot)/q(\cdot)\bigr)$, resampling otherwise.
\item \textbf{Eval harness}: throughput, latency, acceptance rate, and
      distribution-preservation metrics.
\end{enumerate}

\paragraph{Architecture and implementation.}
The engine is written in PyTorch and validated first on CPU (numerical correctness
and unit tests), then deployed on the in-lab Ascend 950 device which we already
have running an open-source inference stack; the speculative loop itself is
independent of vendor kernels. We deliberately \emph{do not} use the speculative
decoding built into vLLM or similar frameworks; the baseline is the same engine
with drafting disabled. A command-line runner and a lightweight web demo expose a
toggle between vanilla and speculative decoding with live tokens/s and acceptance
stats.

\section{Evaluation Plan}
\textbf{Setup.} Target models: Qwen3-8B (primary) and Qwen2.5-7B (generality).
Drafts: Qwen3-0.6B, Qwen2.5-0.5B, and n-gram prompt lookup. A fixed prompt set of
20--50 samples across domains (writing, code, news).

\textbf{Metrics.}
\begin{enumerate}[label=\arabic*.]
\item \emph{Speed:} end-to-end speedup vs.~baseline, tokens/s, time-to-first-token.
\item \emph{Acceptance:} accepted/proposed ratio, drafts "all-correct" fraction.
\item \emph{Quality preservation (the key claim):} greedy outputs must match
      vanilla \emph{token-by-token}; sampling outputs measured by KL divergence /
      perplexity to the vanilla distribution.
\end{enumerate}

\textbf{Ablations.} Draft length $K \in \{1,\dots,8\}$; draft strength (0.6B vs.~0.5B
vs.~n-gram); temperature; batch size.

\section{Expected Outcomes}
\begin{enumerate}[label=\arabic*.]
\item A clean, self-implemented, reproducible speculative-decoding engine with unit
      tests and documented code.
\item Quantitative lossless speedups with systematic ablations and acceptance-rate
      analysis, reported against an honest baseline.
\item A runnable interactive demo (web), the 8-page ACL-style final report, slides,
      demo video, and a source archive with README.
\end{enumerate}

\section{Project Timeline}
\begin{table}[h]
\centering
\small
\begin{tabular}{llp{5.6cm}}
\toprule
\textbf{Phase} & \textbf{Dates} & \textbf{Deliverables} \\
\midrule
Phase 1 & 2026/10/12 & Project proposal (this document, 2pp) \\
Phase 2 & 10/12--10/31 & Engine implemented; first speedup measured on CPU/NPU \\
Phase 3 & 11/01--11/20 & Systematic experiments and ablations; figures \\
Phase 4 & 11/21--11/30 & Interactive demo; midterm report (2--3pp) + slides \\
Phase 5 & 12/01--final & 8-page ACL final report, final slides, demo video, code zip \\
\bottomrule
\end{tabular}
\end{table}

\section{Team Members and Responsibilities}
\begin{table}[h]
\centering
\small
\begin{tabular}{llp{6.2cm}}
\toprule
\textbf{Member} & \textbf{Role} & \textbf{Responsibilities} \\
\midrule
Zou Bowen (2312251) & Developer (Lead) & Core implementation of the
speculative-decoding engine and verification loop (Phase 2); experiment design,
coding and data analysis \\
Cui Di (2412439) & Reviewer, writer \& presenter & Literature collection and
reading (related work); final 8-page ACL-style report; presentation slides and
demo presentation \\
\bottomrule
\end{tabular}
\end{table}

\bibliography{custom}

\end{document}